\section{A Hamiltonian involution of a quadric}
\label{sec:quadric}

We now construct the manifold and involution to which
Proposition~\ref{prop:reduction} will apply.  Let $M\subset\CP^4$ be the
quadric defined in~\eqref{eq:quadric}.

\begin{proposition}\label{prop:quadric}
The manifold $M$ is closed, connected, and of real dimension six. For
the restricted Fubini--Study form $\omega$, normalized below, the map
\[
 A[z_0:z_1:z_2:z_3:z_4]
   =[z_0:-z_1:-z_2:-z_3:-z_4]
\]
is a Hamiltonian involution.  Its fixed set is the smooth submanifold
\[
 F=M\cap\{z_0=0\}\cong S^2\times S^2.
\]
For every $q\in F$, the differential $dA_q$ acts as $+\id$ on $T_qF$
and as $-\id$ on a complementary real two-dimensional subspace of
$T_qM$.
\end{proposition}

\begin{proof}
In any standard affine chart of $\CP^4$, the equation defining $M$
takes the form $1+\sum_j w_j^2=0$.  Its differential cannot vanish on
its zero set: that would force all the $w_j$ to be zero.  Thus $M$ is
a smooth complex hypersurface, of complex dimension three.  As a
closed subset of projective space it is compact, and it has no
boundary.

To see connectedness, write a nonzero isotropic vector as $z=x+iy$
with $x,y\in\R^5$.  Its real and imaginary parts satisfy
\[
 |x|=|y|>0,\qquad x\cdot y=0.
\]
After rescaling, they form an ordered orthonormal pair.  Every such
pair extends to a positively oriented orthonormal basis, so these
pairs form one orbit of the connected group $\mathrm{SO}(5)$.  Their
map $(x,y)\mapsto[x+iy]$ onto $M$ proves that $M$ is connected.

We construct the symplectic form with its normalization explicit.
On the unit sphere $S^9\subset\C^5$, put
\[
 \alpha=\sum_{j=0}^4(x_j\,dy_j-y_j\,dx_j).
\]
The form $d\alpha$ is invariant under scalar phase rotations.  If
$v\in T_zS^9$, then
\[
 d\alpha(iz,v)
   =-2\operatorname{Re}\sum_{j=0}^4\overline{z_j}v_j=0.
\]
Consequently it descends under the quotient map
$p:S^9\to\CP^4$ to a two-form $\sigma$, characterized by
$p^*\sigma=d\alpha$.  Local sections show that $\sigma$ is smooth
and closed.  A projective tangent vector has a horizontal lift
$v\in\C^5$ satisfying $\sum_j\overline{z_j}v_j=0$; multiplication by
$i$ on these lifts induces the complex structure.  Since
\[
 d\alpha(v,iv)=2|v|^2,
\]
$\sigma$ is positive on every complex tangent line.  Its restriction
$\omega=\sigma|_M$ is therefore closed and nondegenerate.

We use a family of Hamiltonian rotations: equal speeds will give the
involution, and unequal speeds will later give a function with four
critical points on $M$. For $a,b,t\in\R$, let $\Phi_t^{a,b}$ keep $z_0$
fixed and act on $(z_1,z_2)^{\mathsf T}$ and $(z_3,z_4)^{\mathsf T}$ by
$R_{at}$ and $R_{bt}$, respectively, where
\[
 R_\theta=
 \begin{pmatrix}
  \cos\theta&-\sin\theta\\
  \sin\theta&\cos\theta
 \end{pmatrix}.
\]
These real orthogonal, complex unitary transformations preserve the
quadric and $\alpha$. If $Z_{a,b}$ is their generating vector field
on $S^9$, Cartan's formula
gives
\[
 \iota_{Z_{a,b}}d\alpha=-d\bigl(\alpha(Z_{a,b})\bigr).
\]
The function $\alpha(Z_{a,b})$ is phase invariant, because the rotations
commute with scalar multiplication.  It thus descends to a smooth
function $g_{a,b}$ on projective space. The displayed identity descends
as well, so $G_{a,b}=-g_{a,b}|_M$ generates the projected rotations on
$M$ with our convention $\iota_{X_{G_{a,b}}}\omega=dG_{a,b}$.
Explicitly,
\begin{equation}\label{eq:rotation-hamiltonian}
 G_{a,b}([z])=
 \frac{2a\,\operatorname{Im}(\overline{z_1}z_2)
       +2b\,\operatorname{Im}(\overline{z_3}z_4)}
      {\sum_{j=0}^4|z_j|^2}.
\end{equation}
In particular, $A=\Phi_1^{\pi,\pi}$ is Hamiltonian and preserves $\omega$.

A projective line fixed by $A$ lies in one of the two eigenspaces of
$\operatorname{diag}(1,-1,-1,-1,-1)$.  The positive eigenspace gives
the single point $[1:0:0:0:0]$, which does not lie in $M$.  The
negative eigenspace gives precisely $F=M\cap\{z_0=0\}$.  To identify
$F$, use the invertible complex-linear map
\[
 (z_1,z_2,z_3,z_4)\longmapsto
 \begin{pmatrix}
  z_1+iz_2&z_3+iz_4\\
  -z_3+iz_4&z_1-iz_2
 \end{pmatrix}.
\]
Its determinant is $z_1^2+z_2^2+z_3^2+z_4^2$.  It therefore
identifies $F$ with the projectivized nonzero rank-one matrices.
The map
\[
 \CP^1\times\CP^1\longrightarrow F,
 \qquad ([u],[v])\longmapsto[uv^{\mathsf T}]
\]
is a diffeomorphism: a nonzero column and row recover the two
projective factors, smoothly on the corresponding matrix charts.
Finally, stereographic coordinates identify each $\CP^1$ with
$S^2$.  Only this smooth identification is needed below.

For the tangent-space assertion, fix $q\in F$ and choose an affine
chart $z_j=1$ with $j>0$.  In its normalized coordinates, $A$ sends
$w_0$ to $-w_0$ and leaves all other coordinates fixed.  The equation
of $M$ is
\[
 w_0^2+1+\sum_{\substack{k>0\\k\ne j}}w_k^2=0.
\]
At $q$, where $w_0=0$, its linearization places no constraint on
the $w_0$ direction.  The remaining differential is nonzero, so its
kernel in the other coordinates is exactly $T_qF$.  Thus
\[
 T_qM=T_qF\oplus\C\frac{\partial}{\partial w_0},
\]
with the asserted $+1$ and $-1$ actions.  This also verifies directly
that $F$ is a smooth complex surface. The positivity established above
therefore makes $\omega|_F$ symplectic.
\end{proof}
